\paragraph{RQ3: feature comparison.}
Table~\ref{tab:ton-features} compares all eight descriptions. With Shin as the SVM baseline, our categories lose 22 correct predictions; world-joint angles gain 67, and both additions together lose three. The same-coordinate angle control gains only four. The larger world-angle gain therefore cannot be attributed to the angle formula alone. The small SVM improvement does not recur with LightGBM.

\begin{table*}[t]\centering\small
\caption{RQ3: pooled classification on 60,251 accepted original images. D: distances; S: directions; E: own categories; J: world-joint measurements. P is macro-precision (\%). All variants contain D.}
\label{tab:ton-features}
\begin{tabular}{lrrrr}\toprule
 & \multicolumn{2}{c}{SVM} & \multicolumn{2}{c}{LightGBM}\\
Hand description & P & F1 & P & F1\\\midrule
\tonTableInput ton_feature_rows_en.tex
\bottomrule\end{tabular}
\end{table*}

In the models with D/S/E, shuffling thumb flexion produces the largest mean F1 drop among our categories for SVM (0.08 percentage points); for LightGBM it is middle--ring spreading (0.03). This identifies model dependence, rather than a guaranteed benefit from adding a feature.

\paragraph{Brief practical explorations.}
Alongside the feature study, we briefly tried ideas for left/right handling and G/H camera views. Mirror training and right normalization help on artificial mirror tests but lose 26 and 136 correct predictions on originals. G/H rotation learns the offered artificial geometry, while two selected live G examples remain H. These small exploratory checks do not establish better recognition for new users; protocols, scores and examples are in Appendix~\ref{app:ton-orientation}--\ref{app:ton-live}.
